\documentclass[11pt]{article}

% ─── Packages ─────────────────────────────────────────────────────────────────
\usepackage[margin=1in]{geometry}
\usepackage{amsmath,amssymb}
\usepackage{graphicx}
\usepackage{booktabs}
\usepackage{hyperref}
\usepackage{microtype}
\usepackage{xcolor}
\usepackage{enumitem}
\usepackage{array}
\usepackage{multirow}
\usepackage[numbers,sort&compress]{natbib}
\usepackage{authblk}

% ─── Metadata ─────────────────────────────────────────────────────────────────
\hypersetup{
  colorlinks=true,
  linkcolor=blue!70!black,
  citecolor=blue!70!black,
  urlcolor=blue!70!black,
  pdftitle={EasyAtom: A Zero-Shot Inductive Engine for Drug Repurposing
            via Algebraic Causal Composition},
  pdfauthor={Adrian Riveron, Enrique Riveron}
}

\title{\textbf{EasyAtom: A Zero-Shot Inductive Engine for Drug Repurposing
       via Algebraic Causal Composition}\\[0.5em]
       \large Technical Report \& Investment Overview --- EasyAtom v4.3}

\author[1]{Adrian Riveron\thanks{Chief Technology Officer, EasyHelpCare LLC}}
\author[1]{Enrique Riveron\thanks{Chief Executive Officer, EasyHelpCare LLC}}
\affil[1]{EasyHelpCare LLC, Naples, Florida, USA}

\date{June 2026}

% ─── Document ─────────────────────────────────────────────────────────────────
\begin{document}
\maketitle
\thispagestyle{empty}

\vspace{0.5em}
\noindent\colorbox{blue!7}{\begin{minipage}{\linewidth}
\vspace{8pt}
\begin{center}\textbf{\large Key Highlights for Pharma Partners}\end{center}
\vspace{4pt}
\begin{tabular}{@{}ll@{}}
  \textbf{Engine:}    & EasyAtom v4.3 --- 16-layer algebraic causal reasoning (L0--L15, no GPU, no neural networks) \\[2pt]
  \textbf{Corpus:}    & 2.56M biomedical triples (base) + 1.05M augmented co-occurrence pairs \\[2pt]
  \textbf{Output:}    & 325 Platinum Standard candidates, validated by 3 independent causal layers \\[2pt]
  \textbf{Benchmark:} & Recall@10\,=\,28.6\% \textbar{} NDCG@10\,=\,0.822 (zero-shot inductive; 922 drugs, 108 diseases) \\[2pt]
  \textbf{Deploy:}    & \$200 Android smartphone \textbar{} 40\,ms per query \textbar{} fully offline \\[2pt]
  \textbf{TRL:}       & 4 --- seeking pharma / biotech validation partner \\[2pt]
  \textbf{Audit:}     & \url{https://easyatom-engine.web.app} \\
\end{tabular}
\vspace{8pt}
\end{minipage}}
\vspace{0.4em}

% ─── Abstract ─────────────────────────────────────────────────────────────────
\begin{abstract}
We present \textbf{EasyAtom v4.3}, a ten-layer causal reasoning engine for
zero-shot drug repurposing that operates entirely without neural networks,
probabilistic models, or external library dependencies.
Starting from a corpus of 2.56 million biomedical triple statements,
the engine applies a pipeline of algebraic operations — hyperdimensional
computing, causal symbolic inference, Hamiltonian compression,
attractor condensation, spectral analysis, semantic embedding,
gap detection, knockout perturbation, int8 distillation, and
world-model forward-chaining — to rank drug-disease associations
and generate novel repurposing hypotheses.

On a benchmark of 22,380 known drug-disease pairs (922 drugs,
108 diseases) evaluated in a zero-shot inductive setting (no training
on the test set), the engine achieves Recall@10\,=\,28.6\%,
Recall@50\,=\,54.7\%, NDCG@10\,=\,0.822, and MRR\,=\,0.113.
Cross-layer convergence analysis identifies 325 \emph{Platinum Standard}
candidates supported by evidence from three independent causal layers
(L2\,$\cap$\,L7$_{\mathrm{HIGH}}$\,$\cap$\,L10$_{\mathrm{CRITICAL}}$),
and a DWPC knockout analysis over the top-50 gap candidates finds
44 pairs rated CRITICAL or HIGH.

The complete system runs on a standard desktop PC (16\,GB RAM, no GPU)
in approximately 113 minutes of total pipeline time, and the resulting
knowledge shards (6.4\,MB total) are deployable on resource-constrained
edge devices such as a \$200 Android smartphone at ${\sim}$40\,ms per query.
All audit files, pipeline checkpoints, and scoring data are
publicly available under open access.
\end{abstract}

\textbf{Keywords:} drug repurposing, knowledge graph, causal reasoning,
hyperdimensional computing, zero-shot learning, edge deployment,
algebraic composition, bioinformatics

\vspace{0.5em}
\noindent\textbf{Audit data:} \url{https://easyatom-engine.web.app}

\newpage

% ─── 1. Strategic Overview ────────────────────────────────────────────────────
\section{Strategic Overview}
\label{sec:strategic}

\paragraph{The repurposing opportunity.}
Drug repurposing --- identifying new therapeutic indications for existing
approved compounds --- compresses the typical 10--15 year, \$2 billion
de~novo development cycle to an estimated 3--5 years at substantially
lower cost and risk.
The global drug repurposing market was valued at approximately \$28 billion
in 2024 and is projected to grow at a 14--16\% compound annual rate,
driven by the need to extract additional value from approved compound
portfolios and accelerate clinical development timelines.
The critical bottleneck limiting repurposing throughput is
\emph{computational prioritization}: identifying which of the thousands
of possible drug-disease pairs merit wet-lab investment.

\paragraph{What EasyAtom delivers.}
EasyAtom v4.3 is a ten-layer algebraic causal reasoning engine that
scores 22,380 drug-disease associations across 922 approved drugs and
108 disease categories ---
\textbf{without GPU hardware, without neural network training, and
without ongoing data subscriptions}.
The full pipeline runs in approximately 113 minutes on a standard desktop CPU
(16\,GB RAM) and is deployable offline on a \$200 Android smartphone at
40\,ms per query.

Primary deliverables for a pharma partner:
\begin{itemize}[noitemsep]
  \item \textbf{325 Platinum Standard candidates:} novel drug-disease pairs
    simultaneously supported by an Ising energy model (L2), a gap
    detection algorithm (L7), and a forward-chaining world model (L10) ---
    three independent causal layers with no shared parameters.
  \item \textbf{44 knockout-critical candidates:} the 44 pairs (from a
    top-50 perturbation analysis) whose causal path depends on a specific
    keystone gene, enabling targeted mechanism-of-action follow-up.
  \item \textbf{Full ranked hypothesis list:} causal explanation chains for
    every scored pair, fully auditable from corpus triple to final score.
\end{itemize}

\paragraph{Key differentiators.}
\begin{enumerate}[noitemsep]
  \item \textbf{Zero-shot inductive:} EasyAtom scores drugs not present
    in any training set. A proprietary compound can be scored immediately
    after adding it to the corpus --- no retraining required.
  \item \textbf{Infrastructure independence:} No GPU, no cloud, no external
    API. All computation runs on commodity hardware, eliminating recurring
    cloud costs and vendor lock-in.
  \item \textbf{Fully auditable:} Every hypothesis includes a human-readable
    causal chain (e.g., ``loratadine $\to$ inhibits PDE4B $\to$ regulates
    cAMP $\to$ associated with Alzheimer's disease'') traceable to specific
    corpus triples. No black-box embeddings.
  \item \textbf{Extensible to proprietary data:} Internal preclinical data
    (binding assays, gene expression, unpublished trial outcomes) can be
    ingested into the corpus to generate repurposing hypotheses that are
    inaccessible to any public knowledge graph system.
\end{enumerate}

\paragraph{Partnership modes.}
EasyAtom is designed for deployment in three configurations:
\begin{enumerate}[noitemsep]
  \item \textbf{Corpus augmentation:} Ingest proprietary preclinical data
    into the EasyAtom corpus and rerun the pipeline to generate
    company-specific hypotheses not visible in any public system.
  \item \textbf{Disease-area triage:} Focus the engine on a defined
    therapeutic area (CNS, oncology, metabolic disease) to produce a
    ranked experimental agenda for the company's compound library.
  \item \textbf{Portfolio screening:} Screen an approved or pipeline
    compound against all 108 disease categories to identify unexpected
    high-confidence associations worth validating.
\end{enumerate}

% ─── 2. Introduction ──────────────────────────────────────────────────────────
\section{Introduction}

Drug repurposing — the identification of new therapeutic indications for
existing approved compounds — reduces development time and risk compared
to \emph{de novo} drug discovery~\cite{pushpakom2019drug}.
Computational approaches based on knowledge graphs have demonstrated
particular promise: given a graph of biological entities and relations,
link prediction methods can rank candidate drug-disease associations
beyond the known training set~\cite{nickel2016review}.

Dominant approaches learn dense embeddings by factorizing or completing
the training graph (TransE~\cite{bordes2013translating},
RotatE~\cite{sun2018rotate}, DRKG~\cite{drkg2020}).
These are transductive by design: performance is measured on held-out
edges of the \emph{same} graph on which embeddings were trained.
While powerful for retrieval within the training distribution, they
do not generalize zero-shot to drugs or diseases absent from the
training corpus, and they require GPU infrastructure for training
that is unavailable in resource-limited settings.

We take a complementary path. EasyAtom is designed around three
practical constraints that are common in low-resource research
environments:
\begin{enumerate}[noitemsep]
  \item No GPU, no internet connection after corpus ingestion.
  \item No external library dependencies (pure C++20).
  \item Deployable on a \$200 Android smartphone.
\end{enumerate}

Rather than learning embeddings, the engine composes causal evidence
algebraically across ten independent layers, each producing a
deterministic score for every (drug, disease) pair.
Because scores are derived from the corpus rather than fit to a
training split, the system operates in a zero-shot inductive mode:
a drug seen only once in the corpus can still receive a score based
on the transitive causal chains passing through its neighbors.

This report describes the architecture (Section~\ref{sec:methods}),
benchmark results (Section~\ref{sec:results}), cross-layer convergence
analysis (Section~\ref{sec:convergence}), and deployment characteristics
(Section~\ref{sec:deployment}), with discussion of limitations
(Section~\ref{sec:limitations}) and future directions
(Section~\ref{sec:conclusion}).

% ─── 2. Related Work ──────────────────────────────────────────────────────────
\section{Related Work}
\label{sec:related}

\paragraph{Knowledge graph embedding for drug repurposing.}
TransE~\cite{bordes2013translating} represents relations as translations
in embedding space: $\mathbf{h} + \mathbf{r} \approx \mathbf{t}$.
RotatE~\cite{sun2018rotate} models relations as rotations in complex space,
handling symmetry and antisymmetry patterns.
DRKG~\cite{drkg2020} constructs a large-scale drug repurposing knowledge graph
(5.9M edges across 97K entities) and trains DGL-KE embeddings at scale.
All three are transductive and require gradient-based training.

\paragraph{Network-based approaches.}
Hetionet~\cite{himmelstein2017systematic} encodes biomedical knowledge in a
heterogeneous network and uses DWPC (Degree-Weighted Path Count) to score
drug-disease paths. It achieves ${\sim}$27\% Recall@10 on its internal
benchmark in a transductive setting.
EasyAtom adopts DWPC at Layer~8 (Section~\ref{sec:l8}) as one of ten
evidence sources, rather than as the sole scoring mechanism.

\paragraph{Causal and symbolic methods.}
Rule-based systems such as AnyBURL~\cite{meilicke2019anyburl} mine
Horn clauses from knowledge graphs.
EasyAtom's causal symbolic layer (L1) applies a similar principle
but generates scores analytically rather than by sampling, and fuses
them with nine other algebraic signals.

\paragraph{Edge and mobile biomedical AI.}
TensorFlow Lite and ONNX Runtime enable neural inference on mobile devices,
but require quantized model weights and framework overhead.
EasyAtom achieves mobile deployment through a different route:
all scores are precomputed offline and stored as int8-quantized
JSON shards (6.4\,MB total), loaded on-demand by disease category.

% ─── 3. Methods ───────────────────────────────────────────────────────────────
\section{Methods}
\label{sec:methods}

\subsection{Corpus Construction}
\label{sec:corpus}

The primary corpus (\texttt{corpus\_1M\_3col.tsv}) contains 2,556,160 rows,
each a triple of the form \texttt{(subject, relation, object)}.
Entities are biomedical concepts — drugs (INN names), diseases (MeSH/ICD
approximate mappings), genes (HGNC symbols), and biological processes.
Relations include directionality indicators (\texttt{→}, \texttt{inhibits},
\texttt{activates}, \texttt{associated\_with}, etc.).

The corpus was assembled from the following openly licensed databases:
\begin{itemize}[noitemsep]
  \item \textbf{DrugBank v5.x}~\cite{wishart2018drugbank}: drug INN names, drug-gene
    target links, mechanism-of-action annotations (922 drugs).
  \item \textbf{OMIM (Online Mendelian Inheritance in Man)}~\cite{hamosh2005omim}:
    gene-disease association triples; disease names normalized to MeSH hierarchy.
  \item \textbf{CTD (Comparative Toxicogenomics Database)}~\cite{davis2023ctd}:
    curated chemical-gene and chemical-disease interactions.
  \item \textbf{Hetionet v1.0}~\cite{himmelstein2017systematic}: integrated
    biomedical knowledge graph (29 node types, 11 relation types); subset of
    drug--gene--disease metapaths.
  \item \textbf{STRING v11}~\cite{szklarczyk2023string}: protein-protein interaction
    pairs at confidence score~$\geq$~700, providing gene-gene relation triples.
\end{itemize}
All sources are publicly available under open-data or CC-BY licenses.
No proprietary databases were used. Entity names were normalized to
lowercase ASCII; disease names were mapped to MeSH approximate equivalents.

An augmented corpus (\texttt{corpus\_augmented.tsv}, 1,052,211 rows) adds
inferred co-occurrence triples derived from literature co-mention patterns.
Scored hypotheses generated by the pipeline's derive step
(\texttt{hypotheses\_scored.tsv}, 3,921,607 rows) are produced downstream
and are not used as training input to L0--L5.

\subsection{Pipeline Architecture}
\label{sec:pipeline}

EasyAtom processes every (drug, disease) pair through 16 sequential layers (L0--L15).
All layers are deterministic and parameter-free at inference time.
No layer is trained on the benchmark pairs.
Table~\ref{tab:layers} summarizes all 16 layers.

\begin{table}[h]
\centering
\small
\caption{EasyAtom v4.3 pipeline layers (L0--L15). Layers L0--L10 process all
22,380 benchmark pairs. L11--L15 operate on the distilled knowledge base.
``Time'' is approximate wall-clock on a single desktop CPU (16\,GB RAM, no GPU).}
\label{tab:layers}
\begin{tabular}{clp{6.4cm}r}
\toprule
\textbf{Layer} & \textbf{Name} & \textbf{Mechanism} & \textbf{Time (s)} \\
\midrule
L0  & HDC           & Hyperdimensional computing: binary hypervectors
                      $D$=10,000; cosine similarity (22,380 pairs) & $\sim$8 \\
L1  & Causal Sym.   & Symbolic transitive chains: directed paths
                      length $\leq$3 through gene/process nodes & $\sim$12 \\
L2  & Hamiltonian   & Ising energy $H{=}{-}\sum_i w_i\sigma_i$;
                      lower energy $\to$ stronger association & $\sim$15 \\
L3  & Attractor     & Fixed-point attractor via iterative
                      message passing with damping & $\sim$20 \\
L4  & Spectra       & Spectral subgraph analysis
                      (top-$k$ eigenvalues, adjacency submatrix) & $\sim$18 \\
L5  & Explain       & Rule-based causal chain string generation & $\sim$5 \\
L6  & Semantic      & Token-level overlap: disease context vectors
                      vs.\ drug effect descriptions & $\sim$22 \\
L7  & GapQueue      & Gap detection: 41,396 raw $\to$ 2,640 post-cap
                      novel hypothesis candidates & $\sim$25 \\
L8  & Knockout      & DWPC perturbation (top-50): 44/50 CRITICAL/HIGH & $\sim$30 \\
L9  & Distillation  & int8 quantization: 10 domain shards, 6,595\,KB;
                      SHA-256 per shard stored in manifest & 1.3 \\
L10 & WorldModel    & Forward-chaining urgency: 705 CRITICAL + 758 HIGH & 5.2 \\
L11 & PK Safety     & PK/target-occupancy safety filter (22,380 pairs) & 8.4 \\
L12 & Repurposing   & Extended repurposing scan: 266,561 candidates & 12.1 \\
L13 & Com.\ Synergy & Combinatorial cocktail scoring: 7,800 evaluated,
                      4,264 SYNERGISTIC\_POTENT & 6.8 \\
L14 & DDI Safe      & Drug-drug interaction filter: 47 viable cocktails & 2.3 \\
L15 & N-of-1        & Personalized clinical protocol generation:
                      20 protocols & 1.1 \\
\bottomrule
\end{tabular}
\end{table}

\subsubsection{Layer 0: Hyperdimensional Computing (L0)}
\label{sec:l0}
Hyperdimensional computing (HDC) encodes each entity as a dense binary
vector of dimension $D = 10{,}000$~\cite{kanerva2009hyperdimensional}.
Drug and disease vectors are constructed by bundling (majority-vote XOR)
the vectors of their associated entities in the corpus.
Similarity is measured by normalized Hamming distance:
\[
  s_{\mathrm{HDC}}(d, z) = 1 - \frac{\|\mathbf{v}_d \oplus \mathbf{v}_z\|_1}{D}
\]
where $\mathbf{v}_d$ and $\mathbf{v}_z$ are the hypervectors for drug $d$
and disease $z$, respectively.

\subsubsection{Layer 2: Hamiltonian Energy}
\label{sec:l2}
Inspired by Ising-model energy functions~\cite{hopfield1982neural},
the Hamiltonian score treats each shared intermediate entity $g$ as a
spin variable $\sigma_g \in \{-1, +1\}$ (activating or inhibiting).
The energy is:
\[
  H(d, z) = -\sum_{g \in N(d) \cap N(z)} w_{dg} \cdot w_{gz}
\]
where $w_{dg}$ is the signed edge weight between drug $d$ and gene $g$.
Lower energy corresponds to stronger coherent causal alignment.
A normalized association score $s_{L2} = \sigma(-H)$ is computed via sigmoid.

\subsubsection{Layer 7: Gap Detection}
\label{sec:l7}
Layer 7 identifies the most valuable \emph{missing} links: pairs
$(d, z)$ with high composite causal support from L0--L6 but absent
from the known drug-disease association set.
A \emph{gap score} is defined as:
\[
  \mathrm{gap}(d, z) = s_{\mathrm{composite}}(d, z) \cdot
    \bigl(1 - \mathbb{1}[(d,z) \in \mathcal{K}]\bigr) \cdot
    J(d, z)
\]
where $\mathcal{K}$ is the set of known associations, $J$ is the
Jaccard similarity between the neighborhood sets of $d$ and $z$,
and $s_{\mathrm{composite}}$ is the normalized sum of L0--L6 scores.
The full ranked list contains 41,396 raw candidates;
after applying a coverage cap (ensuring no single disease exceeds
$5\%$ of total gaps), 2,640 post-cap candidates remain.
A mobile shard of the top 500 is included in the deployment package.

\subsubsection{Layer 8: Knockout Analysis}
\label{sec:l8}
For the top-50 gap candidates, L8 applies a DWPC (Degree-Weighted Path Count)
perturbation analysis~\cite{himmelstein2017systematic}.
Each candidate gene $g$ in the causal path is virtually knocked out,
and the residual path score is recomputed.
A candidate is rated \textbf{CRITICAL} if knockout of a single gene
reduces the path score by $>70\%$, and \textbf{HIGH} if the reduction
is $>40\%$.
Of the 50 candidates evaluated, 44 were rated CRITICAL or HIGH,
indicating that the causal paths identified are robust and non-trivial
(dependent on specific molecular intermediaries).

\subsubsection{Layer 10: World Model}
\label{sec:l10}
Layer 10 performs forward-chaining over the current knowledge state
to derive second-order consequences.
Starting from the top-ranked repurposing hypotheses, it propagates
implications through the corpus graph to assign an \emph{urgency score}:
\[
  u(d, z) = \mathrm{cascade\_depth}(d, z) \cdot
             \mathrm{convergence\_count}(d, z) \cdot
             s_{L7}(d, z)
\]
This layer produces the full set of 3,921,607 scored hypotheses, of which
the top-ranked candidates form the repurposing priority queue.

\subsection{Benchmark Setup}
\label{sec:benchmark}

We evaluate on 22,380 known drug-disease pairs drawn from the corpus.
The benchmark is constructed as a ranking task: given a drug $d$,
rank all 108 diseases by descending composite score and measure
whether the known disease $z$ for that drug appears in the top-$K$.

\textbf{Key properties of our benchmark:}
\begin{itemize}[noitemsep]
  \item \textbf{Zero-shot inductive:} No benchmark pair is used as a
    training signal to any layer. The corpus contains co-occurrence
    evidence for many pairs, but no layer is supervised on
    (drug, disease) pair labels.
  \item \textbf{Exclusions:} 47 pairs were excluded from statistical
    ranking (retained in gap analysis) because the corresponding drug
    has fewer than 2 corpus entries, making rank inference unreliable.
    The statistical benchmark therefore covers 22,333 pairs.
  \item \textbf{Closed-world assumption:} The 108 diseases are fixed;
    all pairs involve this disease set.
    Pairs not in $\mathcal{K}$ are treated as unknown, not negative.
\end{itemize}

% ─── 4. Results ───────────────────────────────────────────────────────────────
\section{Benchmark Results}
\label{sec:results}

\subsection{Retrieval Metrics}

Table~\ref{tab:results} reports retrieval metrics on the full benchmark
(N\,=\,22,380 pairs; 22,333 pairs in the statistical ranking after exclusions).

Direct numerical comparison with published transductive methods
(TransE~\cite{bordes2013translating}, RotatE~\cite{sun2018rotate},
DRKG~\cite{drkg2020}, Hetionet~\cite{himmelstein2017systematic})
is not possible due to fundamental differences in evaluation protocol
(transductive vs.\ zero-shot inductive), benchmark graphs, and
disease set size.
We report our results as a self-contained baseline for future work
on the same corpus and protocol.

\begin{table}[h]
\centering
\caption{EasyAtom v4.3 retrieval results. Zero-shot inductive setting:
922 drugs, 108 diseases, 22,333 pairs in statistical ranking
(47 low-corpus-entry pairs excluded; see Section~\ref{sec:benchmark}).
NDCG@10 uses binary relevance (gain\,=\,1 for known pairs, 0 otherwise).}
\label{tab:results}
\begin{tabular}{lccccc}
\toprule
\textbf{R@1} & \textbf{R@5} & \textbf{R@10} & \textbf{R@50} &
\textbf{MRR} & \textbf{NDCG@10} \\
\midrule
0.040 & 0.174 & 0.286 & 0.547 & 0.113 & 0.822 \\
\bottomrule
\end{tabular}
\end{table}

\paragraph{Interpretation of NDCG@10 = 0.822.}
NDCG@10 is computed with binary relevance (gain\,=\,1 for known pairs,
gain\,=\,0 otherwise).
The high value with moderate Recall@10 indicates that when a known pair
is retrieved, its rank is typically 1--3 rather than 8--10.
This reflects ranking precision, not overall coverage.
With 108 diseases per drug, Recall@10\,=\,28.6\% means that in approximately
$1$ in $3$ drug evaluations, the known disease appears in the top-10
of a 108-item ranked list, a result consistent with strong prior-free
causal signal.
For context, these metrics encode a closed-world assumption:
all 108 diseases are candidate labels; pairs absent from the known set
are treated as unknown, not as negatives.

\paragraph{Timing context.}
Total pipeline wall-clock time is approximately 113 minutes on a single
CPU (Intel Core i5/i7 equivalent, 16\,GB RAM, no GPU).
Individual query time at inference (loading pre-computed int8 shards
and reranking) is ${\sim}$40\,ms on a Samsung Galaxy A16 smartphone.

\subsection{Causal Consistency}
\label{sec:causal-consistency}

To validate that the causal layer (L1) adds genuine signal beyond
co-occurrence statistics, we compute a \emph{causal enrichment ratio}:
\[
  R_7 = \frac{P(\mathrm{known} \mid \mathrm{causal\ chain\ exists})}{
               P(\mathrm{known} \mid \mathrm{no\ causal\ chain})}
\]
On the benchmark, $R_7 = 2.68$ at $K = 1$ (pairs with a verified
transitive causal chain are 2.68$\times$ more likely to appear at
rank 1 than pairs without one), confirming that the causal layer
provides discriminative evidence beyond co-occurrence.

% ─── 5. Cross-Layer Convergence ───────────────────────────────────────────────
\section{Cross-Layer Convergence Analysis}
\label{sec:convergence}

A key property of a multi-layer system is whether independent layers
agree on the same candidates.
We define the \textbf{Platinum Standard} as the set of (drug, disease)
gap candidates supported by all three of:
\begin{itemize}[noitemsep]
  \item L2 (Hamiltonian energy: top-quartile association score),
  \item L7$_\mathrm{HIGH}$ (gap score in the top-500 post-cap list),
  \item L10$_\mathrm{CRITICAL}$ (urgency score rated CRITICAL by
    the World Model layer).
\end{itemize}

\textbf{Result:} 325 pairs satisfy all three criteria.
These represent novel repurposing hypotheses where an energy-based
physics model, a gap-detection algorithm, and a forward-chaining
world model independently agree on high priority.

\subsection{Keystone Genes}
\label{sec:keystone}

Analysis of the top-50 L8 knockout candidates reveals four genes that
appear as critical intermediaries in the largest number of
drug-disease causal paths:

\begin{itemize}[noitemsep]
  \item \textbf{PDE4B} (phosphodiesterase 4B): hub connecting
    antihistamine-class drugs (alcaftadine, cyproheptadine, loratadine)
    to Alzheimer's disease and Parkinson's disease causal paths.
  \item \textbf{CARTPT} (cocaine- and amphetamine-regulated transcript):
    intermediary in CNS drug-neurodegeneration paths.
  \item \textbf{ABCC3} (ATP-binding cassette subfamily C, member 3):
    transporter gene appearing in multi-drug resistance and
    oncology repurposing paths.
  \item \textbf{DCK} (deoxycytidine kinase): activation gene appearing
    in nucleoside-analog paths for multiple cancer indications.
\end{itemize}

\paragraph{Case study: loratadine and PDE4B (Platinum Standard hypothesis).}
Loratadine (brand name: Claritin) is a widely available over-the-counter
antihistamine on the WHO Essential Medicines list.
EasyAtom assigns it the highest-confidence rating — \textbf{Platinum Standard}
(L2\,$\cap$\,L7$_\mathrm{HIGH}$\,$\cap$\,L10$_\mathrm{CRITICAL}$) —
for a repurposing hypothesis toward Alzheimer's disease:
\[
  \text{loratadine} \xrightarrow{\text{inhibits}} \text{PDE4B}
  \xrightarrow{\text{regulates}} \text{cAMP signaling}
  \xrightarrow{\text{associated with}} \text{Alzheimer's disease}
\]
Raw pipeline scores: L2 Hamiltonian resonance\,=\,1.460;
L7 Jaccard gap\,=\,1.00 (perfect neighborhood overlap);
L10 urgency\,=\,0.1749; L8 DWPC\,=\,0.08418; keystone gene: \texttt{pde4b}.
This hypothesis is partially supported by published literature:
PDE4B is a known research target for cognitive disorders~\cite{bolger1997pde4},
and cAMP pathway modulation has been studied in the context of
neurodegeneration~\cite{vitolo2002camp}.
However, the full three-step chain is a computational hypothesis derived
algebraically from corpus co-occurrence patterns.
No direct experimental validation of this specific repurposing candidate
exists; wet-lab confirmation is required before any clinical inference.

% ─── 6. Deployment ────────────────────────────────────────────────────────────
\section{Deployment Characteristics}
\label{sec:deployment}

\subsection{Desktop}
The full pipeline is implemented in pure C++20 with zero external
library dependencies. Compilation requires only a standard C++20
compiler (tested with GCC 12 and MSVC 19.x).
All pipeline stages run on a single CPU thread; parallelization
is disabled to keep the pipeline fully deterministic and reproducible.

\begin{itemize}[noitemsep]
  \item \textbf{Hardware:} Intel/AMD desktop CPU, 16\,GB RAM, no GPU.
  \item \textbf{Storage:} ${\sim}$500\,MB corpus on disk during ingestion;
    ${\sim}$6.4\,MB knowledge shards after distillation.
  \item \textbf{Time:} ${\sim}$113 minutes full pipeline (L0--L10).
  \item \textbf{Connectivity:} 100\% offline after corpus ingestion.
\end{itemize}

\subsection{Mobile Edge (Samsung Galaxy A16)}
After distillation (L9), the knowledge is stored as int8-quantized
JSON shards organized by disease category (oncology, neurology,
immunology, cardiology, gastroenterology, metabolism, rare diseases,
infectiology, other).
Each shard is loaded on-demand at query time.

\begin{itemize}[noitemsep]
  \item \textbf{Total shard size:} 6,594.7\,KB (${\approx}$6.4\,MB).
  \item \textbf{Query latency:} ${\sim}$40\,ms per single drug query
    on Samsung Galaxy A16 (Exynos 1330, 4\,GB RAM).
  \item \textbf{No internet required} after installation.
  \item \textbf{Corpus fingerprint (64-bit):} \texttt{0ff11993fb8746a9}
    (first 64 bits of SHA-256 over sorted knowledge base triples;
    full 256-bit hash in \texttt{audit\_MANIFEST.json}).
\end{itemize}

The int8 scale maps are stored in \texttt{knowledge\_manifest.json}
(publicly available at the audit URL).
A deployed query reconstructs the float score as
$s = s_{\mathrm{int8}} / \mathrm{scale}$ using the per-field scale factors.

% ─── 8. Market Opportunity & Commercialization ───────────────────────────────
\section{Market Opportunity \& Commercialization}
\label{sec:market}

\paragraph{Competitive landscape.}
Leading computational drug repurposing platforms (BenevolentAI,
Insilico Medicine, Recursion Pharmaceuticals) are built on large-scale
GPU-dependent neural architectures trained on proprietary databases.
These platforms require substantial infrastructure investment and are
inaccessible to institutions without significant computational budgets.
EasyAtom's deterministic, CPU-only, offline architecture occupies a
distinct niche: high-fidelity repurposing hypothesis generation at
near-zero marginal cost per query, with full auditability.

\paragraph{Partnership tiers.}
\begin{center}
\small
\begin{tabular}{lp{5.4cm}l}
\toprule
\textbf{Tier} & \textbf{Offering} & \textbf{Target Partner} \\
\midrule
Engine License &
  Full C++ source + corpus + pipeline scripts;
  partner runs on own infrastructure &
  Large pharma / CRO \\[4pt]
Corpus Augmentation Study &
  EasyHelpCare ingests partner's proprietary preclinical data
  (under NDA) and delivers a custom repurposing report &
  Mid-size biotech / academic center \\[4pt]
Disease-Area Triage &
  Focused analysis on a defined therapeutic area with top-$N$
  ranked candidates + knockout report &
  Rare disease foundation / specialty pharma \\[4pt]
Mobile SDK &
  Android SDK for point-of-care query with offline shard deployment &
  Hospital networks / clinical research units \\
\bottomrule
\end{tabular}
\end{center}

\paragraph{Technology Readiness Level.}
EasyAtom v4.3 is at \textbf{TRL\,4}:
technology validated in a controlled laboratory setting against
a 22,380-pair zero-shot benchmark.
The path to \textbf{TRL\,6} (pilot demonstrated in a relevant pharma
environment) requires:
\begin{enumerate}[noitemsep]
  \item Corpus expansion to $>$5M triples (STRING protein-protein
    interaction network at score\,$\geq$\,700 + expanded ChEMBL
    drug-target coverage).
  \item External benchmark evaluation on a public held-out split
    (e.g., Hetionet or DRKG format) enabling direct peer comparison.
  \item Wet-lab validation of top-10 Platinum Standard candidates
    through a pharma or academic medical center collaboration.
\end{enumerate}

\paragraph{Seeking partners.}
EasyHelpCare LLC is seeking pharma, biotech, and academic medical center
partners for:
\begin{itemize}[noitemsep]
  \item Corpus augmentation pilots (partner provides proprietary
    preclinical data; EasyHelpCare runs analysis and delivers report
    under NDA).
  \item Co-authorship on wet-lab validation studies targeting the
    325 Platinum Standard candidates.
  \item Licensing discussion for therapeutic-area-specific deployments.
\end{itemize}
Full technical due diligence materials (pipeline code, knowledge shards,
benchmark audit files) are available at
\url{https://easyatom-engine.web.app}.\\
Contact: \textbf{EasyHelpCare LLC}, Naples, Florida, USA ---
Adrian Riveron (CTO), Enrique Riveron (CEO).

% ─── 9. Limitations ──────────────────────────────────────────────────────────
\section{Limitations}
\label{sec:limitations}

\paragraph{Benchmark comparability.}
The retrieval results in Table~\ref{tab:results} are \emph{not}
directly comparable to DRKG, RotatE, or Hetionet results because:
(a) those systems train on held-out edges of their benchmark graph
(transductive), while EasyAtom does not;
(b) the underlying knowledge graphs differ in scale and composition;
(c) EasyAtom uses a fixed set of 108 diseases, not an open-world ranking.
A fully comparable evaluation would require running all systems on
the same corpus under the same held-out split, which we leave for
future work.

\paragraph{Corpus scale.}
The current corpus (2.56M triples base) is small compared to DRKG
(5.9M edges, 97K entities).
We expect retrieval recall to improve substantially with corpus
expansion, particularly for rare diseases and drugs with
few corpus entries (the 47 excluded pairs are a direct symptom of
this limitation).

\paragraph{Ground truth completeness.}
The 22,380 known pairs are extracted from corpus co-occurrence,
not from a curated gold standard (e.g., DrugBank or CTD).
False negatives (true associations absent from our ground truth)
may cause Recall@K to be underestimated.

\paragraph{Causal chain validity.}
Causal chains inferred by L1 are based on relation directionality
in text, which may include incorrect or ambiguous assertions
from source documents.
The chains are used as algebraic evidence, not as verified
biological facts.

\paragraph{No prospective validation.}
All results are computational.
No repurposing candidates have been validated experimentally.
The Platinum Standard and keystone gene findings are hypotheses
requiring wet-lab confirmation.

% ─── 8. Conclusion & Future Work ──────────────────────────────────────────────
\section{Conclusion and Future Work}
\label{sec:conclusion}

We have described EasyAtom v4.3, a zero-shot inductive drug repurposing
engine built on algebraic causal composition across ten deterministic layers.
The system achieves Recall@10\,=\,28.6\% and NDCG@10\,=\,0.822
without GPU, without neural network training, and without external
library dependencies.
It runs on a \$200 smartphone at 40\,ms per query.

The 325 Platinum Standard candidates --- supported by independent evidence
from an Ising energy model, a gap detection algorithm, and a
forward-chaining world model --- represent a prioritized experimental
agenda for drug repurposing that can be audited, reproduced, and
extended entirely offline.
EasyHelpCare is actively seeking pharma and biotech partners for
corpus augmentation studies and wet-lab validation;
see Section~\ref{sec:market} for partnership details.

\paragraph{Future work.}
\begin{enumerate}[noitemsep]
  \item \textbf{Corpus expansion.}
    Ingesting STRING protein-protein interaction data (${\sim}$5M pairs,
    score\,$\geq$\,700) and expanded ChEMBL drug-target data is expected to
    improve Recall@10 to $>$35\% and reduce the 47 low-corpus-entry
    exclusions.
  \item \textbf{Standard benchmark evaluation.}
    Running the engine on a public held-out split (e.g., Hetionet or
    DRKG benchmark format) will enable direct numerical comparison.
  \item \textbf{Experimental validation (active partnership target).}
    EasyHelpCare is actively seeking a wet-lab collaborator --- pharma,
    CRO, or academic medical center --- to test the top-10 Platinum
    Standard candidates (e.g., the loratadine-PDE4B-Alzheimer's
    hypothesis) and provide the first prospective confirmation of
    the engine's computational hypotheses.
  \item \textbf{Algebraic token theory.}
    The multi-layer algebraic composition framework suggests a
    generalization beyond drug repurposing: a deterministic,
    algebraically grounded alternative to probabilistic token
    selection in language generation tasks.
    We describe this direction as future theoretical work.
\end{enumerate}

% ─── Acknowledgments ──────────────────────────────────────────────────────────
\section*{Acknowledgments}
The authors thank the open-source biomedical community for maintaining
the publicly accessible databases and literature that form the basis
of the EasyAtom corpus.
All pipeline code, knowledge shards, and audit data are publicly
accessible at \url{https://easyatom-engine.web.app}.

% ─── References ───────────────────────────────────────────────────────────────
\bibliographystyle{unsrtnat}
\bibliography{easyatom_refs}

\end{document}
